% =====================================================================
% Pasajele determinate complet de masuratori, gata de lipit.
% Scrise 19 august 2026, dupa regulile de scriitura: fara em-dash, fara
% desemnarea lucrarii noastre prin substantivul generic interzis, fara
% cadru de competitie, contributiile la persoana I. Vezi memoria
% `preferinte-de-scriitura-paper`.
%
% Cifrele vin din output/e_rezumat_modele.csv, output/e_adoptie.md,
% output/e_m1.md, output/e_m3.md si output/repere_publicate.md.
% NU se rotunjesc altfel decat aici fara sa se verifice sursa.
% =====================================================================


% =====================================================================
% 1. ABSTRACT, propozitia de rezultat. Bara comparabila ca protocol e
%    TMU-HIT zero-shot uniform, nu G-Scale. Aia e revendicarea tare.
% =====================================================================

We predict lexical complexity in ten languages without using a single
annotated example from the language being evaluated. Averaged over the
ten languages, we reach a Pearson correlation of $0.675$, ahead of the
only published system evaluated under the same zero-label protocol by
$0.065$, and ahead on nine of the ten languages.


% =====================================================================
% 2. RESULTS, paragraful principal. Ablatia e ordonata ca sa se vada ca
%    cele doua contributii se aduna, nu se suprapun.
% =====================================================================

Table~\ref{tab:results} reports Pearson correlation per language,
averaged without weighting across the ten languages, on 5{,}624 test
instances. The configuration we report reaches $0.6749$. Removing the
encoder average costs $0.0057$ and removing the evaluation-split model
choice costs $0.0052$; removing both costs $0.0104$, with a paired
bootstrap interval of $[+0.006, +0.015]$ over resampled contexts. The
two changes are close to additive, which is what we would expect from
their per-language profiles: the encoder average helps most on Sinhala
and Catalan, the model choice helps almost entirely on French.

All intervals are computed by resampling contexts rather than
instances, because seven of the ten languages contain on average three
target words per context and an instance-level bootstrap would be too
narrow exactly where the dependence is strongest.


% =====================================================================
% 3. RESULTS, comparatia cu reperele. Se scrie in AMBELE sensuri.
%    „Batem cel mai bun sistem publicat" NU se scrie.
% =====================================================================

Two published systems are relevant. TMU-HIT run under a uniform
zero-shot protocol reaches $0.610$, and is the only published number
that uses no target-language labels. We are ahead of it by $0.065$,
winning on nine of the ten languages. G-Scale with GPT-4o reaches
$0.669$ and is the highest published average on this test set. We are
ahead of it by $0.006$.

That second margin should not be read as an improvement over the state
of the art. The 95\% interval of our own average is $[0.660, 0.689]$,
so a difference of $0.006$ is a fifth of one interval width, and it
cannot be tested at all, because per-instance predictions for G-Scale
are not available. The protocols also differ: G-Scale conditions on
thirty labelled examples in the target language and on a temperature
fitted across all ten trial sets, while we use none. Being above it is
therefore not a protocol win, just as being below it would not have
been a protocol failure. We include it because it is the highest
published figure and a reader familiar with this data would notice its
absence.

Under the comparison that does hold protocol fixed, the margin is
$0.065$ rather than $0.006$, and seven of the ten per-language figures
we report are the highest published for that language.


% =====================================================================
% 4. APPROACH sau RESULTS. M1. Asta e pasajul care raspunde obiectiei
%    „ati ales medierea dupa ce ati vazut ca selectia cade".
% =====================================================================

Choosing to average the four encoders rather than select among them is
the one binary decision in our pipeline that could have been made after
inspecting test scores, so we removed it from our hands. For each
target language, we compare the two alternatives using only the other
nine languages, through the same leave-one-language-out criterion that
already fixes the fusion weight and the language model. Averaging is
selected on all ten folds, with a minimum margin of $0.0088$. The
configuration we report is therefore the output of the protocol rather
than a choice we made after seeing the result.

The unanimity is not surprising, and it is worth saying why. Averaging
beats selection on nine of the ten languages, so dropping one language
from the criterion cannot reverse the argument of the maximum. What
this establishes is not that a close decision came out correctly, but
that the decision is stable and does not rest on any single language. A
criterion that flipped when one language was withheld would have been
evidence that the choice was noise, which is precisely what we observe
for per-language encoder selection.

One consequence should be stated plainly. On Filipino, averaging is
worse than selection, and the protocol still selects averaging there,
because it is not allowed to look at Filipino. That is a real loss we
pay to keep the procedure clean.


% =====================================================================
% 5. RESULTS sau LIMITATIONS. Cifra raportata NU e cifra maxima.
%    Se scrie, fiindca materialul suplimentar o contine.
% =====================================================================

The figure we report is not the highest we measured. Averaging four
encoders instead of three reaches $0.6760$, which is $0.0011$ above the
configuration we report, with a paired interval of $[-0.0004, +0.0025]$
that includes zero. We report the three-encoder configuration because
the fourth encoder does not produce a measurable gain while accounting
for close to half of the encoder cost, and because dropping it also
removes a 24-layer model from the requirements. We note this rather
than reporting the maximum, since both figures appear in the
supplementary log.


% =====================================================================
% 6. LIMITATIONS. Bifurcatiile. Se declara, nu se ascunde. Formularea
%    vine din output/e_adoptie.md §4.
% =====================================================================

We measured 51 system configurations on the test split, listed in full
in the supplementary material. Each of them respects the
leave-one-language-out protocol at every level of choice, so no
per-language figure uses labels from the language it is evaluated on.
The choice of which configuration to report, however, was made by us,
by looking at the mean test score, and is not itself the output of an
automatic procedure. The figure we report should be read as the
outcome of a search of that size rather than a first attempt. Confidence
intervals are paired against the preceding configuration and are not
corrected for multiple comparisons.

We reduced this exposure where it was cheapest to do so. The single
decision most open to the objection, whether to average or select
encoders, is now taken by the same leave-one-language-out criterion
used elsewhere. The remaining choices are declared here instead.


% =====================================================================
% 7. LIMITATIONS. Restrangerea bazinului. Se scrie pana cand P10 ruleaza;
%    daca P10 ruleaza, paragraful se INLOCUIESTE, nu se sterge.
% =====================================================================

The pool from which the language model is selected is narrower on the
test split than on the trial split. Two of the three candidates were
run on all ten languages, while the third was run only on French, so
the selection effectively chooses between two candidates per language.
The restriction uses no labels and is a matter of coverage rather than
of score, but it does narrow the choice, and a reader could reasonably
observe that we selected among the models we happened to have run.


% =====================================================================
% 8. LIMITATIONS. n, acum o nota de subsol si nu o limitare.
% =====================================================================

We evaluate on 5{,}624 of the 5{,}627 test instances. The three
remaining items are cases where the target string could not be located
in its context by any of the matching rules we apply.


% =====================================================================
% 9. TABELUL DE REZULTATE. Cifrele per limba sunt din
%    output/fusion_lolo_adoptat.csv; reperele din repere_publicate.md,
%    care e transcris de mana din Nohejl et al., Tabelul 13, si din
%    Anexa B a lui Shardlow et al. 2024a pentru cele doua randuri TMU-HIT
%    uniforme, care NU exista publicate ca medii.
% =====================================================================

\begin{table*}[t]
\centering
\small
\begin{tabular}{lccccccccccc}
\toprule
& ca & en & fil & de & fr & it & ja & pt & si & es & avg \\
\midrule
Frequency baseline      & .301 & .748 & .389 & .591 & .517 & .519 & .642 & .713 & -.196 & .551 & .478 \\
Archaeology             & .274 & .790 & .443 & .551 & .533 & .534 & .485 & .683 & .044 & .527 & .487 \\
TMU-HIT, zero-shot      & .528 & .804 & \textbf{.569} & .658 & .625 & .539 & .645 & .764 & .248 & .720 & .610 \\
TMU-HIT, per-language best & \textbf{.616} & .820 & \textbf{.569} & .718 & .625 & .601 & .733 & .786 & .308 & .762 & .654 \\
G-Eval, GPT-4o          & .573 & .827 & .494 & .711 & .764 & .629 & .745 & .761 & .328 & .739 & .657 \\
G-Scale, GPT-4o         & .583 & .840 & .504 & .714 & \textbf{.774} & .654 & .767 & .774 & .330 & .749 & .669 \\
\midrule
Ours                    & .541 & \textbf{.870} & .512 & \textbf{.733} & .645 & \textbf{.665} & \textbf{.784} & \textbf{.843} & \textbf{.396} & \textbf{.761} & \textbf{.675} \\
\bottomrule
\end{tabular}
\caption{Pearson correlation per language on the MLSP test set, and the
unweighted mean over the ten languages. Bold marks the highest figure
in each column. Only the frequency baseline, Archaeology and our own
row use no labels in the language being evaluated; TMU-HIT zero-shot
uses none either, while the remaining rows condition on three to thirty
target-language examples.}
\label{tab:results}
\end{table*}


% =====================================================================
% 10. CONTRIBUTII, pentru introducere. Paragraf narativ, persoana I,
%     fara buline, fara cifre de sprijin, fara formule.
% =====================================================================

First, we show that the choice between selecting a component per
language and averaging over all candidates can be made by the same
leave-one-language-out criterion that fixes every other choice, and
that the criterion gives opposite answers for two different components
of the same pipeline, each confirmed by measurement. Second, we show
that fitting a component on the trial split while fitting the rest on
the evaluation split is a measurable error rather than a detail, and
that correcting it recovers most of the gap on the language where the
gap was largest. Third, we replace character length with tokenizer
fertility and resolve frequency lookups through lemmas, two changes
motivated by morphology that transfer across scripts. Fourth, we report
every configuration we measured, together with what each choice
between them cost us.
